% 本文件由 main.tex 输入，作为“实证结果分析”章节，不单独编译。

\section{实证结果分析}

在前文理论分析基础上，本文进一步利用现有可用结果，对政府产业引导基金促进城市创新的基准效应、债务压力的调节作用及其潜在机制进行系统检验。为保证内容边界清晰，本章仅依据现有可用结果文件展开，不扩展到描述性统计原始表或新增回归设定之外的内容。

\subsection{研究设计}

\subsubsection{样本、变量与估计方法}

研究单元为地级市—年份观测，样本期为2015—2024年。因变量采用城市创新产出，核心口径包括发明申请量、实用新型申请量和专利申请总量。核心解释变量为政府产业引导基金设立口径，其中累计设立规模是全文主线变量，基金当年设立数量和累计设立数量作为辅助口径。调节变量为债务压力及其滞后一期口径，机制变量依次对应早期投资事件占比、社会资本撬动效率以及地区金融发展水平。

控制变量保持与现有回归结果一致，主要包括经济发展水平、财政科技支出、人口规模、产业结构和外资利用水平。所有主回归均采用城市固定效应和年份固定效应，并按城市维度聚类稳健标准误。由于不同变量存在缺失或口径差异，具体样本量随模型而变化，因此后文对每张表均单独报告观测值数量。

\subsubsection{模型设定}

基准回归用于检验政府产业引导基金是否能够促进城市创新，调节模型进一步识别债务压力是否会削弱基金的创新促进作用，机制模型则围绕耐心资本属性、社会资本撬动效率和融资约束缓解效应展开。相应地，可将主模型写为：

\begin{equation}
Innovation_{it}=\alpha+\beta Fund_{it}+\gamma Controls_{it}+\mu_i+\lambda_t+\varepsilon_{it}.
\end{equation}

\begin{equation}
Innovation_{it}=\alpha+\beta_1 Fund_{it}+\beta_2 Debt_{it-1}+\beta_3(Fund_{it}\times Debt_{it-1})+\gamma Controls_{it}+\mu_i+\lambda_t+\varepsilon_{it}.
\end{equation}

\begin{equation}
M_{it}=\alpha+\theta_1(Fund_{it}\times Debt_{it-1})+\gamma Controls_{it}+\mu_i+\lambda_t+\nu_{it},
\end{equation}

\begin{equation}
Innovation_{i,t+k}=\alpha+\rho_1(Fund_{it}\times Debt_{it-1})+\rho_2 M_{it}+\rho_3(Fund_{it}\times M_{it})+\gamma Controls_{it}+\mu_i+\lambda_t+\eta_{it}.
\end{equation}
\ERJWhere{$Innovation_{it}$ 表示城市创新产出，$Fund_{it}$ 表示政府产业引导基金口径，$Debt_{it-1}$ 表示滞后一期债务压力，$M_{it}$ 表示机制变量，$Controls_{it}$ 为一组城市层面控制变量，$\mu_i$ 和 $\lambda_t$ 分别表示城市固定效应和年份固定效应。耐心资本与金融发展机制主要关注 $\rho_2$，社会资本撬动效率机制进一步关注 $\rho_3$。}

\subsection{基准回归结果}

基准回归首先回答政府产业引导基金是否能够提升城市创新。现有结果显示，无论采用基金设立数量还是累计设立规模，政府产业引导基金扩张与城市专利申请产出之间均存在稳定正相关关系。其中，累计设立规模在后续债务调节效应与机制检验中也保持了最强的一致性，因此下文将其作为正文主解释变量。

\begin{table}[htbp]
  \centering
  \caption{政府产业引导基金与城市创新：基准回归主结果}
  {\fangsong\zihao{-5}
  \begin{tabularx}{0.96\linewidth}{l l c c c}
    \toprule
    因变量 & 基金口径 & 系数 & p值 & N \\
    \midrule
    发明申请量 & 累计设立规模 & 0.0300*** & $<0.001$ & 1083 \\
    实用新型申请量 & 累计设立规模 & 0.0405*** & 0.0031 & 1083 \\
    专利申请总量 & 累计设立规模 & 0.0800** & 0.0174 & 1083 \\
    发明申请量 & 当年设立数量 & 297.1917*** & 0.0007 & 1149 \\
    专利申请总量 & 当年设立数量 & 849.5984*** & 0.0018 & 1149 \\
    发明申请量 & 累计设立数量 & 106.7968*** & 0.0007 & 1149 \\
    \bottomrule
  \end{tabularx}

  \par\smallskip
  {\songti\zihao{6}注：所有模型均控制城市固定效应、年份固定效应以及经济发展水平、财政科技支出、人口规模、产业结构和外资利用水平，并按城市聚类稳健标准误。表中星号分别表示在 10\%、5\%、1\% 水平上显著。由于当年基金设立规模口径结果不稳，正文主线不予采用。}
  }
\end{table}

从表1可以看出，累计设立规模对三类创新产出均呈显著正向影响，说明引导基金并非仅仅增加财政性投入规模，而是与城市专利申请增长存在系统性的正向联系。当年设立数量和累计设立数量的估计结果也保持显著为正，这意味着政府产业引导基金的扩张无论从“数量”还是“规模”视角观察，都与城市创新产出改善相一致。结合后续结果的连贯性，本文据此认为假说 H1 获得较强支持。

\subsection{债务压力的调节效应}

在确认基金总体上具有创新促进效应之后，进一步需要识别这一效应是否会受到地方债务约束的影响。现有结果表明，当将债务压力引入模型后，基金累计设立规模与债务压力的交互项在各类创新产出下均显著为负，并且滞后一期债务压力的结果更符合因果时序，因此本文将其作为主规格。

\begin{table}[htbp]
  \centering
  \caption{债务压力对基金创新效应的调节作用}
  {\fangsong\zihao{-5}
  \begin{tabularx}{0.94\linewidth}{l c c c}
    \toprule
    因变量 & 交互项系数 & p值 & N \\
    \midrule
    专利申请总量 & -0.0000609*** & $<0.001$ & 914 \\
    发明申请量 & -0.0000162*** & $<0.001$ & 914 \\
    实用新型申请量 & -0.0000292*** & $<0.001$ & 914 \\
    \bottomrule
  \end{tabularx}

  \par\smallskip
  {\songti\zihao{6}注：表中交互项为“基金累计设立规模$\times$滞后一期债务压力”。所有模型均加入控制变量、城市固定效应和年份固定效应，并按城市聚类稳健标准误。同期债务压力口径的交互项同样显著为负，但本文优先报告滞后一期债务压力结果。财政压力口径在现有结果中呈正向交互，不纳入正文主结论。}
  }
\end{table}

表2显示，在债务压力较高的城市中，同等规模的政府产业引导基金难以产生与低债务压力城市同样强的创新促进效果。无论考察发明申请、实用新型申请还是专利申请总量，这一负向调节效应都具有高度稳定性。由此可以认为，债务压力确实显著削弱了政府产业引导基金促进创新的边际作用，假说 H2 得到有力支持。

\subsection{机制检验}

在主效应和调节效应成立之后，进一步的问题是：债务压力究竟通过何种途径削弱基金的创新促进作用。基于现有可用结果，本文主要从耐心资本属性、社会资本撬动效率和融资约束缓解效应三条路径展开检验。

\subsubsection{耐心资本属性机制}

第一条机制聚焦于基金的耐心资本属性。理论上，地方债务压力越高，越可能促使基金配置偏向短期、低风险项目，进而削弱其“投早、投小、投长期”的特征。现有结果使用早期投资事件占比作为代理变量，并对比例变量的分母稳定性进行了专门处理。为避免单笔投资导致占比机械地落在0或1，样本限定为基金投资事件数不少于2的城市—年份观测。该处理并非装饰性筛选，而是识别比例机制所必需的稳健化步骤。

\begin{table}[htbp]
  \centering
  \caption{耐心资本属性机制：5\% 水平主结果}
  {\fangsong\zihao{-5}
  \resizebox{\linewidth}{!}{%
  \begin{tabular}{l l l c c c c c c}
    \toprule
    机制变量口径 & 创新结果 & 债务处理 & 基准交互项 & 机制方程 & 结果方程 & 加机制后债务调节项系数 & N & $R^2$ \\
    \midrule
    原始比例 & 两年后专利申请总量 & 原始债务压力 & -0.085** & -0.143** & 0.092** & -0.071* & 266 & 0.680 \\
    比例1/99缩尾 & 两年后专利申请总量 & 原始债务压力 & -0.085** & -0.143** & 0.092** & -0.071* & 266 & 0.680 \\
    原始比例 & 两年后专利申请总量 & 债务压力1/99缩尾 & -0.075** & -0.128** & 0.091** & -0.064* & 266 & 0.680 \\
    比例1/99缩尾 & 两年后专利申请总量 & 债务压力1/99缩尾 & -0.075** & -0.128** & 0.091** & -0.064* & 266 & 0.680 \\
    \bottomrule
  \end{tabular}%
  }

  \par\smallskip
  {\songti\zihao{6}注：基准交互项表示未加入机制变量时“基金规模$\times$债务压力”对未来创新的影响；机制方程表示该交互项对早期投资事件占比的影响；结果方程表示早期投资事件占比对未来创新的影响；加机制后债务调节项系数为结果方程中继续保留的债务调节项。表中4条结果均为分母稳定化样本下达到5\% 显著水平的主规格，城市固定效应、年份固定效应和控制变量均已纳入。}
  }
\end{table}

表3显示，在基金投资事件数不少于2的样本中，债务压力显著降低了早期投资事件占比，而早期投资事件占比又显著提升未来两年的专利申请总量。与此同时，基金规模与债务压力交互项的总效应显著为负；加入早期投资事件占比后，该交互项仍为负，但绝对值和显著性均有所下降，表明债务压力不仅直接压缩基金促进创新的边际效应，还会通过削弱早期投资偏好这一耐心资本特征，进一步抑制后续创新产出。需要说明的是，该结论主要依赖分母稳定化样本，因此应理解为透明的有条件证据，而非对全部样本的无条件结论。

\subsubsection{社会资本撬动效率机制}

第二条机制关注社会资本撬动效率。现有结果并未支持一个普遍适用的普通中介链条，而是更倾向于支持一种调节式或承接机制：债务压力先削弱社会资本撬动效率及其改善，较高的撬动效率再强化基金规模对创新的边际促进作用。需要特别强调的是，这一证据主要集中在低金融发展地区，因此其适用范围具有明显的情境依赖。

\begin{table}[htbp]
  \centering
  \caption{社会资本撬动效率机制：最强证据}
  {\fangsong\zihao{-5}
  \resizebox{\linewidth}{!}{%
  \begin{tabular}{l l c c c c c}
    \toprule
    因变量 & 债务处理 & 基准交互项 & 机制方程 & 结果方程 & 加机制后债务调节项系数 & N \\
    \midrule
    实用新型申请量对数 & 滞后债务压力 & -4.03E-08*** & -8.95E-08** & 6.31E-07** & -3.68E-08*** & 705 \\
    实用新型申请量对数 & 滞后债务压力缩尾 & -4.03E-08*** & -8.95E-08** & 6.31E-07** & -3.68E-08*** & 705 \\
    下一期专利申请总量 & 滞后债务压力 & -8.20E-04** & -8.95E-08** & 0.02418** & -6.61E-04*** & 705 \\
    下一期专利申请总量 & 滞后债务压力缩尾 & -8.20E-04** & -8.95E-08** & 0.02418** & -6.61E-04*** & 705 \\
    \bottomrule
  \end{tabular}%
  }

  \par\smallskip
  {\songti\zihao{6}注：表中结果均来自低金融发展、剔除2020年的样本，并采用“债务交互项影响机制变量、基金与机制变量交互影响创新”的方案二扩展模型。机制变量为社会资本撬动效率年度变化指标。加入机制项后，原债务调节项绝对值下降，说明存在承接机制。}
  }
\end{table}

从表4可以看到，在低金融发展地区，债务压力会削弱社会资本撬动效率，而更高的撬动效率能够强化基金规模的创新促进作用。更重要的是，在纳入机制项之后，原先基金与债务压力交互项的绝对值出现下降，说明债务压力对基金创新效应的负向调节，部分是通过社会资本撬动效率这一承接渠道体现出来的。结合现有证据的分布情况，本文将该机制表述为“低金融发展地区中的调节式机制”，而不将其扩展为普遍适用的普通中介结论。

\subsubsection{融资约束缓解机制}

第三条机制检验基金能否通过改善地区金融发展环境来缓解融资约束。本文采用存款/GDP口径金融发展水平作为融资约束的反向代理：存款/GDP越高，地区金融资源供给和吸纳能力越强，融资约束通常越弱。在已有可用结果中，最适合正文的是存款/GDP口径金融发展水平的反双曲正弦变换结果。

\begin{table}[htbp]
  \centering
  \caption{融资约束缓解机制：金融发展口径结果}
  {\fangsong\zihao{-5}
  \resizebox{\linewidth}{!}{%
  \begin{tabular}{l l l c c c c c c}
    \toprule
    机制变量口径 & 创新结果 & 债务处理 & 基准交互项 & 机制方程 & 结果方程 & 加机制后债务调节项系数 & N & $R^2$ \\
    \midrule
    存款/GDP水平（asinh变换） & 专利申请总量对数 & 滞后债务压力对数 & -2.14E-09 & -3.29E-08** & 0.404** & 1.11E-08 & 944 & 0.703 \\
    \bottomrule
  \end{tabular}%
  }

  \par\smallskip
  {\songti\zihao{6}注：基准交互项表示未加入机制变量时“基金累计设立规模$\times \ln(1+\text{滞后一期债务压力})$”对创新的影响；机制方程表示该交互项对存款/GDP金融发展水平的影响；结果方程表示存款/GDP金融发展水平对创新的影响；加机制后债务调节项系数为结果方程中继续保留的债务调节项。机制变量数值越大表示融资约束越弱。该规格纳入控制变量、城市固定效应、年份固定效应，并按城市聚类稳健标准误。}
  }
\end{table}

表5表明，债务压力较高时，政府产业引导基金对存款/GDP口径金融发展水平的改善作用会显著减弱，而金融发展水平提升又会正向促进创新产出。加入机制变量后的债务调节项不显著，因此该表主要提供两步机制链条证据，而不是依赖交互项衰减幅度。这意味着债务压力不仅直接压缩基金的创新促进效应，还会通过削弱基金改善金融环境、缓解融资约束的能力，间接影响城市创新。

\subsection{本节小结}

基于现有可用结果，本文的实证分析可以归纳为以下几点。第一，政府产业引导基金总体上能够提升城市创新产出，尤其是基金累计设立规模与城市专利申请之间存在稳定的正向关系。第二，地方债务压力显著削弱了基金的创新促进效应，这一负向调节在发明申请、实用新型申请和专利申请总量三个口径下均保持稳健。第三，机制检验显示，债务压力可能通过削弱基金投向早期项目的比例、降低社会资本撬动效率以及弱化地区金融发展改善效应三条路径，进一步抑制基金对创新的促进作用。

与此同时，现有机制证据也存在清晰边界。耐心资本机制依赖于基金投资事件数不少于2的分母稳定化样本；社会资本撬动效率机制主要出现在低金融发展地区，属于情境依赖较强的承接机制；金融发展机制虽然在理论链条上较为完整，但其关键显著性处于5\% 边界附近。因而，本文更适合将机制结果理解为与主结论方向一致、但具有样本条件和识别边界的补充证据。总体而言，政府产业引导基金的创新促进作用并非无条件成立，其政策绩效在很大程度上取决于地方财政债务约束与地区资本市场环境。
